%!TeX TS-program = lualatex
%!TeX encoding = UTF-8 Unicode
%!TeX spellcheck = fr-FR
%!BIB TS-program = biber
% -*- coding: UTF-8; -*-
% vim: set fenc=utf-8
\documentclass[10pt,french,a4paper]{article}%,oneside,draft,twoside
%============ common ===================
% \usepackage[utf8]{inputenc}%
% \usepackage[official]{eurosym}
\usepackage[french]{babel}%
\usepackage{csquotes}%
\usepackage{etaremune}%
\usepackage[final]{pdfpages}%
%% --- unit formatting ---
\usepackage{microtype}	% Better typography \tolerance 1414
\hbadness 1414
\emergencystretch 1.5em
\hfuzz 0.3pt
\widowpenalty=10000
\vfuzz \hfuzz
\raggedbottom
%%%%%%%%%%%%%%%%%%%%%%%%%%%%%%%%%%%%%%%%%%%%%%%%%%%%%%%%%%%%%%%%%%%%%
\usepackage{amsmath,amsfonts,amsthm}			% Math packages
\usepackage{bm}%This package defines commands to access bold math symbols
%============ graphics ===================
\graphicspath{{./figures/}}%
\usepackage{graphicx}					% Enable pdflatex
\usepackage{subfig}

%============ front-matter ===================
\newcommand{\BookTitle}{Centre National de la Recherche Scientifique\\de décembre 2022 à juin 2025}
\newcommand{\Title}{%
% Curriculum Vit\ae\ détaillé
%Programme de recherche%: La dynamique de la vision est un processus prédictif
Rapport d'activité scientifique
}%
\newcommand{\SubTitle}{%
Pour évaluation par les sections du Comité national}%
\newcommand{\Author}{Laurent U~Perrinet}%
\newcommand{\Team}{\'Equipe NEural OPerations in TOpographies (NeOpTo)}%
\newcommand{\Institute}{Institut de Neurosciences de la Timone}%
\newcommand{\InstituteUMR}{UMR 7289, CNRS~/~Aix-Marseille Université}%
\newcommand{\Address}{27, Bd. Jean Moulin, 13385 Marseille Cedex 5, France}
\newcommand{\Website}{https://laurentperrinet.github.io/}
\newcommand{\Email}{Laurent.Perrinet@univ-amu.fr}
%%%%%%%%%%%%%%%%%%
%BIBLIOGRAPHY 
%%%%%%%%%%%%%%%%%%
\usepackage[natbib=true,
			bibencoding=utf8,
%			encoding=utf8,
			maxcitenames=7,
			maxnames=7,
			%minnames=3,
			maxbibnames=7,
            sortcites=true,
            block=space,
            backend=biber,
%            doi=false,isbn=false,url=false,
            citestyle=alphabetic,%authoryear-comp,
            backref=true,
%            bibstyle=alphabetic
            ]{biblatex}%
\addbibresource{~/sdrive_cnrs/blog/perrinet_curriculum-vitae.tex/LaurentPerrinet_publications.bib}
\addbibresource{~/sdrive_cnrs/blog/perrinet_curriculum-vitae.tex/LaurentPerrinet_talks.bib}
% \addbibresource{~/Zotero/better-bibtex/zotero_export.bib}
%\usepackage[normalem]{ulem} %for \uline{}
%\DeclareNameFormat{given-family}{
%  \ifgiveninits
%  {\ifthenelse{\equal{\namepartfamily}{Toto}}
%    {\uline{\namepartfamily\addspace\namepartgiveni\namepartsuffix}}
%    {\namepartfamily\addspace\namepartgiveni\namepartsuffix}
%    \ifthenelse{\value{listcount} < \value{liststop}}
%    {\addcomma}
%    {\ifthenelse{\ifmorenames}{~et \,al \adddot}}
%    {}
%  }
%}
% https://tex.stackexchange.com/questions/359311/how-to-underline-name-of-specific-authors-in-biblatex
\usepackage{xpatch}
\usepackage[normalem]{ulem}

\newbox\savenamebox

%\newbibmacro*{name:bbold}[2]{%
%  \def\do##1{\iffieldequalstr{hash}{##1}{\bfseries\setbox\savenamebox\hbox\bgroup \listbreak}{}}%
%  \dolistloop{\boldnames}%
%}
\newbibmacro*{name:bbold}[2]{%
  \def\do##1{\iffieldequalstr{hash}{##1}{\setbox\savenamebox\hbox\bgroup \listbreak}{}}%
  \dolistloop{\boldnames}%
}

\newbibmacro*{name:ebold}[2]{%
  \def\do##1{\iffieldequalstr{hash}{##1}{\egroup\uline{\usebox\savenamebox}\listbreak}{}}%
  \dolistloop{\boldnames}%
}

\xpatchbibmacro{name:given-family}{\usebibmacro{name:delim}{#2#3#1}}{\usebibmacro{name:delim}{#2#3#1}\begingroup\usebibmacro{name:bbold}{#1}{#2}}{}{}

%\xpretobibmacro{name:given-family}{\begingroup\usebibmacro{name:bbold}{#1}{#2}}{}{}
\xapptobibmacro{name:delim}{\begingroup\normalfont}{}{}

\xapptobibmacro{name:given-family}{\usebibmacro{name:ebold}{#1}{#2}\endgroup}{}{}
\xapptobibmacro{name:delim}{\endgroup}{}{}

\newcommand*{\boldnames}{}
\forcsvlist{\listadd\boldnames}{
  {8fa031352aa34db95f5c1022358042a3}, % InBold
  {01b588ba4e4ad753feae6c81709fc04b}} % Highlight

%%%%%%%%%%%%%%%%%%%%%%%%%%%%%%
%% OPTIONAL MACRO FILES
%\usepackage{fltpage}
%\usepackage{tikz}
%%%%%%%%%%%%%%%%%%
%\usepackage{multicol}
%\usepackage{nag}
%============ hyperref ===================
\usepackage[unicode,linkcolor=blue,citecolor=blue,filecolor=black,urlcolor=blue,%pdfborder={0 0 0}
]{hyperref}%
\hypersetup{%
unicode = true, %
pdftitle={\Title},%
pdfauthor={\Author < \Email > - \Institute, \InstituteUMR , \Address - \Website},%
pdfsubject={\Title}%
}%
%\hypersetup{linkcolor=blue,citecolor=blue,filecolor=black,urlcolor=blue}
%\hyphenpenalty=5000
%\tolerance=1000
%% DOCUMENT LAYOUT
%\usepackage{geometry}
%\geometry{a4paper, textwidth=5.5in, textheight=8.5in, marginparsep=7pt, marginparwidth=.6in}
%\setlength\parindent{0in}
%% ---- MARGIN YEARS
\newcommand{\years}[1]{\marginpar{\textit{\scriptsize #1}}}
\providecommand{\natexlab}[1]{#1}
%\providecommand{\url}[1]{\texttt{#1}}
%\expandafter\ifx\csname urlstyle\endcsname\relax
 \providecommand{\doi}[1]{doi: #1}%\else
%  \providecommand{\doi}{doi: \begingroup \urlstyle{rm}\Url}\fi
%\makepagestyle{chapter}
%\makeoddfoot{chapter}{\addRevisionData}{\thepage}{}
%\makeevenfoot{chapter}{\addRevisionData}{\thepage}{}

%%%%%%%%%%%% Her begynner selve dokumentet %%%%%%%%%%%%%%%

% for units
\usepackage{siunitx}%
\newcommand{\ms}{\si{\milli\second}}%
\newcommand{\m}{\si{\meter}}%
\newcommand{\s}{\si{\second}}%
\usepackage{geometry}
\geometry{hscale=0.66,vscale=0.8,centering}

\begin{document}

\begin{titlepage}

\begin{center}
%\vskip 2cm
\emph{\Large \BookTitle }%
\vskip 2cm
{\Huge \Title}
\vskip 1cm
{\Large Laurent U~\textsc{Perrinet}}
\vskip 1cm
\emph{\Large \SubTitle }\\%
%\begin{tabular}[t]{ccc}
%\includegraphics[height=2.5cm]{logo_cnrs-UAM.png} &%\includegraphics[]{}
%\includegraphics[height=2.5cm]{logo_int.png} %& %
%%	\includegraphics[height=2.5cm]{U2_logo.pdf}
%\end{tabular}
\vskip 1cm
\includegraphics[width=1.\textwidth]{~/sdrive_cnrs/blog/perrinet_curriculum-vitae.tex/logotypes/troislogos.png}

\vskip 1cm
\begin{tabular}[t]{|c|}
\hline
\Team \\\hline
\Institute \\ \InstituteUMR \\\hline
\Address \\\hline
\url{\Website}\\\hline
\url{\Email}\\\hline
\end{tabular}
\vskip .5cm
\vfill
{\large 11 Mars 2026}
\end{center}
\end{titlepage}

\newpage

Ce document synthétise mon activité de chercheur CNRS au sein de l'Institut de Neurosciences de la Timone (UMR 7289) durant la période Décembre 2022 à Juin 2025. Il synthétise mes avancées en neurosciences computationnelles combinant expérimentation, modélisation et applications, notamment dans le champ de l'Intelligence Artificielle (IA). Parmi les résultats clés~: la découverte de neurones résilients à l'incertitude dans le cortex visuel (\emph{Nature Communications Biology}, 2023), le développement d'algorithmes bio-inspirés pour des robots aériens (ANR AgileNeuRobot), et une théorie innovante du codage temporel neuronal (projet Polychronies). Ces travaux ont donné lieu dans cette période à 14~publications, 5~thèses encadrées, et des collaborations internationales, tout en s'inscrivant dans une démarche active de vulgarisation (\emph{The Conversation}, \emph{Cerveau \& Psycho}). Ce rapport est écrit pour une évaluation à mi-vague par la section d'évaluation du CoNRS au printemps 2026 et suit dans sa forme les recommandations du comité.


\tableofcontents
\newpage
%%%%%%%%%%%% Endyet daag koepf dokumentet %%%%%%%%%%%%%%%

\section{Curriculum Vit\ae\ }
\subsection{Présentation rapide}%
%\subsection{ Curriculum Vit\ae\ : Laurent U Perrinet}%
%\begin{resume}
\indent 53 ans, né le 23 Février 1973 à Talence (Gironde, France).  %\\

\begin{itemize}
	\item Directeur de Recherche (DR2, CNRS), Affiliation: \Team\ - \Institute\ (\InstituteUMR)
	\item Adresse: \Address
	\item E-mail: \url{mailto:\Email}
	\item Téléphone: 04 91 32 40 44
	\item URL: \url{\Website}
\end{itemize}

\subsection*{Bio-sketch}
Laurent Perrinet est chercheur en neurosciences au CNRS, à l'Institut de Neurosciences de la Timone à Marseille. Il s'intéresse à une question fondamentale : comment notre cerveau parvient-il à construire une image cohérente du monde à partir d'informations visuelles fragmentées et incertaines ?
Pour y répondre, il développe des modèles informatiques inspirés du fonctionnement du cerveau. Ces modèles reproduisent la façon dont les neurones communiquent par de brèves impulsions électriques, et cherchent à expliquer pourquoi la vision biologique est à la fois si rapide et si efficace énergétiquement. Les algorithmes qui en résultent ouvrent des pistes concrètes pour une intelligence artificielle plus sobre et plus robuste.

\subsection*{Résumé quantitatif des contributions}

\begin{itemize}
	
	\item	$63$ publications dans des revues avec comité de lecture dont $4$ en préparation ou en révision, $16$ en premier auteur, $24$ en dernier auteur, avec un total de $4231$ citations, indice h$=30$ (h$=20$ depuis 2021) et indice i10$=61$\footnote{Au 8 Mars 2026, cf.~\url{https://scholar.google.com/citations?user=TVyUV38AAAAJ&hl=fr}},
	\item	$148$ publications dans des actes de congrès avec comité de lecture,
	\item	$4$ livres et $6$ chapitres de livres,
	\item	$77$ conférences invitées avec présentations orales
	ou affichées (dont $21$ dans des congrès internationaux),
	\item \href{https://www.theses.fr/074493701}{Thèses de doctorat}:
	\begin{itemize}
		\item $2$ en cours en encadrement principal (Alexandre Lainé, Matthis Dallain),
		\item $3$ en cours co-dirigées (Agathe Choplin, Charlotte Roy, Kevin Mairot),
		\item $5$ finalisées en encadrement principal~\citep{Khoei14thesis,Boutin20thesis,Franciosini21thesis,Ladret24thesis,Grimaldi24thesis,Jeremie25thesis} soit $3$ dans la période de référence,
		\item $4$ co-dirigées finalisées~\citep{Kremkow09thesis,Damasse18thesis,MansourPour19thesis,Gruel23thesis},
		\item rapporteur de $18$ thèses ($7$ dans la période de référence), Président du jury de jury pour $4$ thèses ($3$), Examinateur pour $13$ thèses ($8$).
	\end{itemize} %
	\item $5$ Post-docs dirigés (Wahiba Taouali, Nicole Voges, Alberto Vergani, Adrien Fois, Thomas Kronland-Martinet),
	\item $1$ contrat en PI local ($150$\,k€, contrat total $435$\,k€), 1 AAP ($100$\,k€), $3$ bourses de thèse obtenues,
	\item Participation à $20$ contrats en collaborateur (dont $12$ ANRs).
	
	
\end{itemize} %

\subsection*{Mots clés}
Perception, vision, détection du mouvement. Calcul parallèle événementiel, émergence dans les systèmes complexes, codage neural, réseaux de neurones impulsionnels. Inférence Bayesienne, minimisation de l'énergie libre, statistiques des scènes naturelles.

\newpage 
\subsection{Diplômes \& titres universitaires}

%\begin{itemize}
\textbf{Habilitation à Diriger des Recherches}, AMU, Marseille\hfill \years{\textbf{2017}}\\
\vspace*{-.15in}
\begin{itemize}
\item[] \'Ecole Doctorale Sciences de la Vie et de la Santé, Aix-Marseille Université, France.
\item[] Sous le titre ``Codage prédictif dans les transformations visuo-motrices'', j'ai défendu mon Habilitation à Diriger des Recherches le 14 avril 2017.
% \item[] Le jury était constitué des Prof. Laurent Madelain (Université Lille III), Dr. Alain Destexhe (Université Paris XI, Rapporteur), Prof. Gustavo Deco (Universitat Pompeu Fabra, Barcelona, Rapporteur), Dr. Guillaume Masson (Aix-Marseille Université), Dr. Viktor Jirsa (Aix-Marseille Université, Rapporteur) et du Prof. Jean-Louis Mege (Aix-Marseille Université).
\end{itemize} %

%%%%%%%%%%%%%%%%%%%%%%%%%%%% These %%%%%%%%%%%%%%%%%%%%%%%%%%%%%%%%%%%%%%%%%%%
\vspace*{.3cm}
	%\item[]
	\textbf{Doctorat de Sciences cognitives} ONERA/DTIM, Toulouse \hfill \years{\textbf{1999-2003}} \\
\vspace*{-.15in}
\begin{itemize}
\item[] Titre : \emph{Comment déchiffrer le code impulsionnel de la Vision? \'Etude du flux parallèle, asynchrone et épars  dans le traitement visuel ultra-rapide.}  Allocataire d'une bourse MENRT, accueil à  l'ONERA/DTIM.
% \begin{itemize}
%\vspace*{.05in}
		\item % codirection
		Cette thèse a été initiée par les résultats de la collaboration pendant le stage de DEA. Elle a été dirigée par Manuel Samuelides (professeur à  \textsc{Supaéro} et chargé de recherche à  l'ONERA/DTIM) et co-dirigée par Simon Thorpe (directeur de recherche au \textsc{CerCo})
		% \item Participation et présentations à  de nombreux colloques internationaux (IJCNN99, NeuroColt00, CNS00, CNS01, LFTNC01, ESANN02, NSI02). Participation aux écoles d'été ``EU Advanced Course in Computational Neuroscience'' à  Trieste (Italie) et ``Telluride Neuromorphic Workshop'' au Colorado (\'Etats-Unis).
		% \item % + org. dynn / enseignement (supaero trex pir / ensica/ cea) / chap. dynn
		% En parallèle, j'ai participé à  l'organisation d'une conférence sur les réseaux de neurones dynamiques (\textsc{Dynn}*2000). Je me suis aussi impliqué dans des activités d'enseignement (initiation à  la programmation sous Matlab et théorie de la probabilité) pour des élèves de première et deuxième année d'école d'ingénieur (à  \textsc{Supaéro} et à  l'\textsc{Ensica}, Toulouse) et des travaux dirigés de robotique (Traitement de l'image et reconnaissance d'objets au laboratoire d'Informatique et d'Automatique de \textsc{Supaéro}). %Enfin, je me suis impliqué dans la rédaction d'un chapitre d'un livre
		\item % soutenance
		La thèse de doctorat a été soutenue le 7 février 2003 à  l'université Paul Sabatier avec la mention "Très honorable avec les félicitations du jury". Le jury était présidé par Michel Imbert (Prof. Université P. Sabatier) et constitué par Yves Burnod (Directeur de recherche à  l'\textsc{INSERM} U483) et Jeanny Hérault (Professeur à  l'INPG, Grenoble).
	% \end{itemize} %
\end{itemize} %


%%%%%%%%%%%%%%%%%%%%%%%%%%%% dea %%%%%%%%%%%%%%%%%%%%%%%%%%%%%%%%%%%%%%%%%%%
\vspace*{.3cm}
	%\item[]
	\textbf{DEA de Sciences cognitives} \hfill \years{\textbf{1998-1999}} \\ Univ. Paris VII, P. Sabatier, EHESS, Polytechnique, mention TB. Allocataire d'une bourse de DEA.
	%%  (stage fin étude simon / dea paris - petitot / stage cert / 2?/26, mention bien)
	% \begin{itemize}
	% 		\item  Assistant de recherche, ONERA/DTIM \years{3/1999-7/1999}(Département de Traitement de l'Image et de Modélisation), Toulouse (stage de DEA).
	% 		\begin{itemize}
	% 			\item  \'Etude de l'apprentissage de type Hebbien de réseaux de neurones basés sur un codage par rang.
	% 			\item  Application à  la reconnaissance de textures visuelles.
	% 		\end{itemize} %
	% 		\item Assistant de recherche, USAFB (Rome, NY)  \years{7/1999-8/1999}/ University of San Diego in California (\'Etats-Unis). % https://www.rollingstone.com/music/music-news/19-worst-things-about-woodstock-99-176052/
	% 		%\begin{itemize}\item
	% 		\'Etude de l'apprentissage autonome dans un système complexe de type automate cellulaire. %\item Application à  la modélisation du comportement d'un pilote d'avion.\end{itemize} %
	% \end{itemize} %

%%%%%%%%%%%%%%%%%%%%%%%%%%%% supaero %%%%%%%%%%%%%%%%%%%%%%%%%%%%%%%%%%%%%%%%%%%
\vspace*{.3cm}
	%\item[]
	\textbf{Diplôme d'ingénieur {\sc Supaéro}}, Toulouse, France.  \hfill \years{\textbf{1993-1998}}\\
Spécialisation dans le traitement du signal et de l'image et en particulier dans les techniques des réseaux de neurones artificiels. \\
% \vspace*{-.1in}
% 	\begin{itemize}
% 	 	\item Projets individuels sur la perception visuelle, la reconnaissance de locuteur, la compression de la parole et sur la reconnaissance de caractères.
% 		\item Ingénieur \textsc{Alcatel}, Vienne (Autriche). \years{9/1995-6/1996} Département du \emph{Voice Processing Systems}. Ce 'stage long' volontaire, intégré à  une formation de \textsc{Supaéro} sur les systèmes industriels, impliquait toutes les étapes de la conception d'un système de messagerie téléphonique de technologie élevée : conception, prototype, contrôle de qualité et étude marketing. %Le stage s'intégrait.
% 		\item  Assistant de recherche, \textsc{Jet Propulsion Laboratory}  \years{4/1997-9/1997}(\textsc{Nasa}), Pasadena, Californie. Département des Sciences de la Terre, Laboratoire d'imagerie radar, \textsc{Interférométrie radar SAR appliquée à la géophysique}
% 		\begin{itemize}
% 			\item Programmation d'un processus de traitement d'images radar interférométriques SAR comprenant des corrections géographiques, une série de filtrages et un traitement d'interférométrie.
% 			\item \'Etude et programmation d'un générateur de carte topographique.% (DEM).
% 			\item Traitement des images obtenues pour surveiller la déformation de la croûte terrestre. \'Etude des tremblements de terre de Landers (Californie) et de Gulan (Chine).
% 		\end{itemize}
% 		\item Assistant de recherche, \textsc{CerCo}  \years{4/1998-9/1998}(CNRS, UMR5549), Toulouse (stage de fin d'études d'ingénieur). Développement d'un réseau de neurones asynchrone appliqué à  la reconnaissance de caractères.
% 		\begin{itemize}
% 			\item Programmation du code du réseau de neurones asynchrones.
% 			\item \'Etude et utilisation des statistiques non-paramétriques %\citep{Barbe95}
% pour l'utili\-sation d'un code basé sur le rang d'activation des neurones. \item Implantation d'une nouvelle règle d'apprentissage du réseau de neurones.
% 		\end{itemize} %
% 	\end{itemize}

\subsection{Expérience scientifique professionnelle}
%%%%%%%%%% CNRS %%%%%%%%%%%%%%%%%%%%%%%%%%
\vspace*{.1cm}
	\textbf{Directeur de Recherche (DR2)} (section CID51), INT/CNRS, Marseille  \years{\textbf{2020-\ldots}} \\


	\textbf{Chargé de Recherche Classe Normale}, INT/CNRS, Marseille  \years{\textbf{2019-2020}} \\
	 % (\emph{sous réserve de financement})
	 Au 1er janvier 2019, j'ai intégré l'équipe \href{http://www.int.univ-amu.fr/spip.php?page=equipe&equipe=NeOpTo&lang=en}{NeOpTo} de Frédéric Chavane (DR, CNRS). J'implémente les modèles prédictifs dans des architectures bio-mimétiques.


	\textbf{Chargé de Recherche grade 1}, INT/CNRS, Marseille  \years{\textbf{2012-2019}} \\
	 Au 1er janvier 2012, notre équipe a intégré l'\Institute\ (\InstituteUMR\ ) à  Marseille (direction Guillaume Masson). J'ai approfondi les modèles en me concentrant sur un codage probabiliste distribué appliqué à la boucle sensori-motrice.


	%\item[]
\textbf{Mission longue} \years{10/2010-01/2012}  Karl Friston's theoretical neurobiology group (The Wellcome Trust Centre for Neuroimaging, University College London, London, UK). Collaboration avec Karl Friston sur l'application de modèles d'énergie libre aux mouvements oculaires.

	\textbf{Chargé de Recherche grade 2} (section 7), INCM/CNRS, Marseille  \years{\textbf{2004-2012}} \\
	 Sous la conduite de Guillaume Masson à  l'INCM à  Marseille, j'ai étudié des modèles spatio-temporels d'inférence dans des scènes naturelles en application de la compréhension des mouvements oculaires.
%%%%%%%%%%%%%%%%%%%%%%%%%%%% post-doc %%%%%%%%%%%%%%%%%%%%%%%%%%%%%%%%%%%%%%%%%%%

%\vspace*{.3cm}
	%\item[]
	\textbf{Post-doctorat}, Redwood Neuroscience Institute (RNI), \'Etats-Unis  \years{\textbf{2004}} \\	% San Francisco (
	 Sous la conduite de Bruno Olshausen, j'ai comparé des modèles standards d'apprentissage avec une méthode originale centrée sur les potentiels d'action. Notamment, j'ai développé une méthode générique évaluant l'importance des processus homéostatiques dans l'apprentissage non-supervisé, qui a conduit à une publication dans le journal Neural Computation (référence A20-\citep{Perrinet10shl}).   J'ai ensuite étendu ce modèle à  l'apprentissage spatio-temporels dans des flux video.


%\vspace{.25in}
\newpage


\section{Résumé de mon activité scientifique}

Mon programme de recherche vise à élucider les relations entre structure
neurale et fonction sensorielle : comment l'organisation topographique des
neurones et la nature impulsionnelle du signal nerveux permettent-elles aux
systèmes sensoriels de s'adapter de façon optimale aux statistiques des
scènes naturelles, par des mécanismes de traitement prédictif ?
Plus précisément, je cherche à formaliser les facultés sensorielles et
cognitives sous la forme de réseaux de neurones impulsionnels réalisant
efficacement des algorithmes de perception visuelle. Les potentiels d'action
--- brèves impulsions se propageant de neurone en neurone --- constituent
une caractéristique universelle des systèmes nerveux et offrent un substrat
naturel pour modéliser le traitement dynamique de l'information. Sur le plan
fonctionnel, j'intègre dans ces modèles des stratégies d'inférence reposant
sur des mécanismes d'apprentissage auto-organisés qui ancrent les relations
spatio-temporelles entre neurones. Sur le plan applicatif, ces travaux
ouvrent la voie à de nouveaux algorithmes neuromorphiques à haute efficacité
énergétique.

\subsection*{Travaux antérieurs}

Mes travaux de thèse, dirigés par Simon Thorpe et Manuel Samuelides, ont
posé les bases mathématiques d'un codage neural impulsionnel pour les images
statiques~\citep{Perrinet03ieee}, en montrant notamment comment des
représentations parcimonieuses peuvent émerger de règles d'apprentissage
hebbiennes~\citep{Perrinet10shl}. En collaboration avec Guillaume Masson,
ces modèles ont été étendus à l'inférence statistique appliquée aux
mouvements des yeux et à la boucle perception--action~\citep{Simoncini12}.
La collaboration avec Karl Friston a ensuite permis de fonder un cadre
théorique rigoureux pour intégrer les délais temporels dans les modèles
prédictifs~\citep{PerrinetAdamsFriston14}.

\subsection*{Avancées récentes}

Ces fondations ont guidé une série d'extensions vers le traitement dynamique
de scènes visuelles en mouvement~\citep{KhoeiMassonPerrinet17}, puis vers
des architectures hiérarchiques bio-inspirées intégrant parcimonie et
rétroaction descendante~\citep{BoutinFranciosiniChavaneRuffierPerrinet20,Franciosini21}.
En collaboration avec Frédéric Chavane, nous avons revisité la structure de
la connectivité horizontale dans le cortex visuel primaire~\citep{Chavane22}
et mis en évidence le rôle de la précision sensorielle dans le traitement
cortical de l'information~\citep{Ladret23}, ouvrant la voie à une
représentation parcimonieuse plus efficace des images naturelles~\citep{Ladret24sparse}.

\subsection*{Perspectives}

Ces avancées ont permis de concevoir des algorithmes neuromorphiques de type
réseaux de neurones impulsionnels (SNN) capables d'une reconnaissance robuste
et ultra-rapide fondée sur la détection de motifs spatio-temporels de
spikes~\citep{Grimaldi24}. L'intégration d'une rétinotopie fovéale dans des
architectures bio-inspirées améliore par ailleurs significativement la
localisation et la catégorisation d'objets dans des conditions
écologiques~\citep{Jeremie25,Jeremie25thesis}. Des collaborations récentes
ont enfin étendu le cadre du codage prédictif à la prévision de séries
temporelles~\citep{Gunasekaran25} et à la modélisation de textures dynamiques
naturelles~\citep{Meso25}, illustrant la portée transversale de ces approches
bien au-delà de la vision.


\newpage
\section{Rapport d’activité sur la période de décembre 2022 à juin 2025} %{\small(journaux à  comité de lecture)}}

% \input{activité}

\newpage
\subsection{Productions scientifiques sélectionnées}

Je présente ici un choix de 10 productions scientifiques pour les 10 derniers
semestres d’activité. Les codes dans la marge correspondent à ceux utilisés dans la \href{https://raw.githubusercontent.com/laurentperrinet/perrinet_curriculum-vitae.tex/master/perrinet_publications.pdf}{liste complète des publications}.

% Publications sélectionnées de la période Décembre 2022 -- Juin 2025
% Suivant le modèle du rapport d'activité CNRS
% Numérotation continuant depuis A50 (dernier numéro connu)

\begin{enumerate}

% -------------------------------------------------------
% REVUES INTERNATIONALES À COMITÉ DE LECTURE
% -------------------------------------------------------

\item[A58] \fullcite{Meso25}% _LA_
\begin{itemize}
  \item Cette publication présente une boîte à outils open source pour la génération en temps réel de textures dynamiques naturelles (Motion Clouds), permettant leur intégration directe dans des environnements expérimentaux de psychophysique (Psychtoolbox, PsychoPy). Elle fournit ainsi un outil standard pour les études de perception visuelle du mouvement.
\end{itemize}

\item[A57] \fullcite{Gunasekaran25}% _LA_
\begin{itemize}
  \item Nous étendons le cadre du codage prédictif à la prévision de séries temporelles. Un mécanisme de rétroaction dynamique, inspiré des principes de minimisation de la surprise, permet au modèle de s'adapter aux dérives de distribution des données. Appliqué à la prédiction de crises épileptiques, le gain est de l'ordre de 40\,\% en AUC-ROC.
\end{itemize}

\item[A56] \fullcite{Ladret24sparse}% _LA_
\begin{itemize}
  \item Ce travail montre que l'introduction d'une hétérogénéité dans les noyaux de filtrage améliore la parcimonie des représentations d'images naturelles. Ce résultat relie directement la diversité observée dans les champs récepteurs biologiques à une fonction computationnelle précise : maximiser l'efficacité du codage.
\end{itemize}

\item[A55] \fullcite{Grimaldi24}% _LA_
\begin{itemize}
  \item Nous proposons une approche événementielle robuste pour la reconnaissance d'objets en continu (\emph{always-on}), à partir de flux de caméras événementielles. L'originalité est de combiner un réseau de neurones à impulsions (SNN) avec une caractérisation extensive des performances en fonction des paramètres d'entrée, en rapprochant les modèles neuromorphiques de l'efficacité du système biologique.
\end{itemize}

\item[A54] \fullcite{Grimaldi23BC}% _LA_
\begin{itemize}
  \item Cette publication introduit une couche de neurones à impulsions capable d'apprendre des délais hétérogènes pour la détection ultra-rapide du mouvement. Un résultat clé est un gain énergétique de l'ordre de 500 fois par rapport aux approches analogiques classiques, tout en maintenant des performances comparables sur des données neuromorphiques réelles.
\end{itemize}

\item[A53] \fullcite{Ladret23}% _LA_
\begin{itemize}
  \item Ce travail expérimental et computationnel révèle deux populations neuronales distinctes dans le cortex visuel primaire (V1) : l'une encodant l'orientation avec précision, l'autre exhibant une invariance à la variance sensorielle via des boucles récurrentes. Ces résultats remettent en question les modèles classiques de traitement visuel hiérarchique et ouvrent la voie à des algorithmes robustes au bruit.
\end{itemize}

\item[A52] \fullcite{Franciosini21}% _LA_
\begin{itemize}
  \item Ce travail étend un modèle prédictif de V1 en y intégrant un mécanisme de \emph{pooling} bio-inspiré, permettant d'expliquer à la fois la diversité fonctionnelle et structurelle observée entre espèces dans le cortex visuel primaire. Il établit un lien quantitatif entre les paramètres architecturaux du modèle et les données anatomiques.
\end{itemize}

\item[A51] \fullcite{Chavane22}% _LA_
\begin{itemize}
  \item Dans cette revue de l'état de l'art sur la connectivité horizontale de V1, nous proposons une hypothèse novatrice : la connectivité s'organise non pas selon le principe classique \og like-to-like \fg{} (neurones de même orientation), mais de façon plus distribuée, \og like-to-all \fg{}. Cette hypothèse réconcilie des données anatomiques et fonctionnelles jusqu'alors contradictoires.
\end{itemize}

% -------------------------------------------------------
% ACTES DE CONGRÈS AVEC COMITÉ DE LECTURE
% -------------------------------------------------------

\item[C9] \fullcite{Nunes23ICCV}%
\begin{itemize}
  \item Nous présentons un algorithme temps réel pour le calcul de cartes de temps-de-contact à partir d'une seule caméra événementielle, par estimation conjointe de la profondeur inverse et du mouvement global. Cette contribution est essentielle pour la navigation autonome de robots aériens à faible consommation.
\end{itemize}

\item[C8] \fullcite{Grimaldi22polychronies}%
\begin{itemize}
  \item À notre initiative, nous développons les fondements d'une théorie du codage neural basée sur des motifs spatio-temporels précis d'impulsions (\emph{polychronies}). Nous passons en revue les évidences biologiques et neuromorphiques en faveur de ce paradigme, et proposons des méthodes de détection adaptées.
\end{itemize}


\end{enumerate}
\newpage
\section{Enseignement, formation et diffusion de la culture scientifique de décembre 2022 à juin 2025} %{\small(journaux à  comité de lecture)}}

\subsection{Encadrement de thèse et post-doctorants} %

Cette section résume les encadrements de \href{https://www.theses.fr/074493701}{thèses de doctorat}:

Actuellement, j'ai l'occasion d'encadrer cinq doctorants %et un post doctorant
en tant qu'encadrant principal ou en co-supervision:

\begin{itemize}
	\item \href{https://laurentperrinet.github.io/author/alexandre-lain%C3%A9/}{Alexandre Lainé} ``Model-based analysis of neurobiological data'' (PhD, bourse école doctorale, 10/2024-09/2027, en co-direction avec Sophie Deneve, INT)
	\item \href{https://laurentperrinet.github.io/author/matthis-dallain/}{Matthis Dallain} ``Focus of attention: a sensory-motor task for energy reduction in spiking neural networks.'' (PhD, 10/2024-09/2027, en co-direction avec Benoît Miramond et Laurent Rodriguez, LEAT)
	\item \href{https://laurentperrinet.github.io/author/kevin-mairot/}{Kevin Mairot} ``L’intelligence artificielle comme aide au diagnostic des dystrophies rétiniennes'' (PhD, bourse de médecine, 10/2023-09/2027)
	\item En co-supervision : \href{https://laurentperrinet.github.io/author/agathe-choplin/}{Agathe Choplin} ``Characterization of Physiological States using Machine Learning'' (PhD, 10/2024-09/2027)
	\item En co-supervision : \href{https://laurentperrinet.github.io/author/charlotte-roy/}{Charlotte Roy} ``Impact de l'incarnation dans un avatar sur le bien-être et la prise de décision dans le Métaverse'' (PhD, 10/2024-09/2027, en co-supervision avec Olivia Petit, KEDGE)
\end{itemize}


Précédemment, j'ai eu l'occasion d'encadrer des étudiants en direction de thèse, en post-doctorat ou en co-direction de thèse :
\begin{itemize}
	\item \href{https://laurentperrinet.github.io/author/jean-nicolas-jeremie/}{Jean-Nicolas Jérémie} ``Foveal Retinotopy and Dual Pathways: A Computational Model for Active Visual Search''~\textcite{Jeremie25thesis} (PhD, bourse AgileNeuRobot (ANR-20-CE23-0021), 10/2021-09/2025)
	\item \href{https://laurentperrinet.github.io/author/thomas-kronland-martinet/}{Thomas Kronland-Martinet} ``Ultra-fast processing of sensor inputs acquired by an event camera'' (Post-doc, bourse AgileNeuRobot (ANR-20-CE23-0021), 01/2025-07/2025)
	\item \href{https://laurentperrinet.github.io/author/adrien-fois/}{Adrien Fois} ``Accurate detection of precise spiking motifs in neurobiological data'' (Post-doc, bourse Polychronies grant (AMX-21-RID-025), 09/2023-03/2025)
	\item \href{https://laurentperrinet.github.io/author/antoine-grimaldi/}{Antoine Grimaldi} ``Ultra-fast vision using Spiking Neural Networks'' (PhD, APROVIS3D grant~\textcite{Grimaldi24thesis} (ANR-19-CHR3-0008-03), 2020 / 2023)
	\item \href{https://laurentperrinet.github.io/author/amélie-gruel/}{Amélie Gruel} ``Design of bio-inspired spiking neural networks (spiking neurons) for event-based stereovision'' (PhD, APROVIS3D grant~\textcite{Gruel23thesis} (ANR-19-CHR3-0008-03), 2020 / 2023)
	\item \href{https://laurentperrinet.github.io/author/hugo-ladret/}{Hugo Ladret} ``A multiscale cortical model to account for orientation selectivity in natural-like stimulations''~\textcite{Ladret24thesis}  (direction, bourse AMU, co-direction avec Christian Casanova, en cotutelle avec l'Université de Montréal, 2019 / 2023)
	% \item \href{https://laurentperrinet.github.io/author/alberto-arturo-vergani/}{Alberto Arturo Vergani} 	``Visual computations using Spatio-temporal Diffusion Kernels and Traveling Waves'' (Post-Doc, 04/2020 - 09/2021)
	% \item \href{https://laurentperrinet.github.io/author/victor-boutin/}{Victor Boutin} 	``Controlling an aerial robot by human semaphore gestures using a bio-inspired neural network''~\textcite{Boutin20thesis} (PhD, bourse AMIDEX, 12/2016-03/2020)
	% \item \href{https://laurentperrinet.github.io/author/angelo-franciosini/}{Angelo Franciosini} ``Trajectories in natural images and the sensory processing of contours''~\textcite{Franciosini21thesis} (PhD, bourse PhD program, 2017 / 2021)
	% \item \href{https://laurentperrinet.github.io/author/kiana-mansour-pour/}{Kiana Mansour Pour} 	``Predicting and selecting sensory events: inference for smooth eye movements''~\citep{MansourPour19thesis} (PhD, 2015 - 2018)
	% \item \href{https://laurentperrinet.github.io/author/jean-bernard-damasse/}{Jean-Bernard Damasse} 	``Smooth pursuit eye movements and learning: Role of motion probability and reinforcement contingencies''~\citep{Damasse18thesis} (PhD, 2014-2017)
	% \item \href{https://laurentperrinet.github.io/author/mina-a-khoei/}{Mina A Khoei} 	``Emerging properties in a neural field model implementing probabilistic prediction''~\textcite{Khoei14thesis} (PhD, 2011-2014)
	% \item \href{https://laurentperrinet.github.io/author/wahiba-taouali/}{Wahiba Taouali} 	``Motion Integration By V1 Population'' (Post-Doc, 2013-2015)
	% \item \href{https://laurentperrinet.github.io/author/nicole-voges/}{Nicole Voges} 	``Complex dynamics in recurrent cortical networks based on spatially realistic connectivities'' (Post-Doc, 2008-2010)
	% \item \href{https://laurentperrinet.github.io/author/jens-kremkow/}{Jens Kremkow} ``Correlating Excitation and Inhibition in Visual Cortical Circuits: Functional Consequences and Biological Feasibility''~\citep{Kremkow09thesis} (PhD, 2006-2009)
%	\item David Arbib 	OBV1: Sélectivité à l'orientation dans le cortex visuel primaire. (master student, 2016-01 / 2016-06)
%	\item Jean Spezia 	developing MotionClouds (undergrad, 2014-09 / 2014-12)
\end{itemize}


\subsection{Participation à des jurys de thèse dans la période de référence} %

\begin{itemize}

    % \item ``Amélioration de l'apprentissage dans les réseaux de neurones impulsionnels'' par Ilyass Hammouamri sous la direction de Timothée Masquelier au CERCO - Centre de Recherche Cerveau et Cognition (Toulouse), soutenue en Octobre 2025,

    % \item ``Apprentissage supervisé efficace pour les réseaux de neurones impulsionnels à codage temporel'' par Gaspard Goupy sous la direction de Ioan Marius Bilasco et Pierre Tirilly  dans l'équipe FOX du laboratoire CRIStAL (Lille), soutenue en Octobre 2025,

    \item ``Méthodes d'apprentissage profond inspirées par la physique pour l'inférence de la qualité de l'air'' par Matthieu Dabrowski sous la direction de Chaabane Djeraba, Jérôme Riedi et José Mennesson - Informatique et applications - Université de Lille, Soutenue en 2024,

    \item ``Avancées en vision neuromorphique : représentation événementielle, réseaux de neurones impulsionnels supervisés et pré-entraînement auto-supervisé'' par Sami Barchid sous la direction de Chaabane Djeraba - Informatique et applications - Université de Lille, Soutenue en 2023,

    \item ``Co-simulation de modèles de circuits neuronaux à macro- et méso-échelle : une technologie en avance sur son temps'' par Lionel Kusch sous la direction de Viktor K. Jirsa --- Neurosciences, Aix-Marseille Université, soutenue en 2023 (Président du jury),

    \item ``Autonomous learning in a neuromorphic vision system : an active efficient coding spiking model applied to vision'' par Thomas Barbier sous la direction de Jochen Triesch - Informatique - Université Clermont Auvergne, Soutenue en 2023  (Examinateur),

    \item ``Approches neuro-robotique intégrées pour la localisation et la navigation d'un véhicule autonome'' par Sylvain Colomer sous la direction de Olivier Romain - STIC (Sciences et Technologies de l'Information et de la Communication) - ED EM2PSI - CY Cergy Paris Université, Soutenue en 2022,

    \item ``Theoretical framework for Time-To-First-Spike coding in Spiking Neural Networks'' par Lina del Pilar Bonilla Camelo sous la direction de Timothée Masquelier - Image, Information, Hypermédia - Toulouse 3, Soutenue en 2022 (Président du jury),

	\item ``Codage prédictif dans le cerveau et les réseaux de neurones profonds'' 	par Zhaoyang Pang sous la direction de Rufin VanRullen - Toulouse 3 (Neurosciences), Soutenue le 28-02-2022,

    \item ``Deep learning and neuroscience : a match made in heaven?'' par Bhavin Yogesh Choksi sous la direction de Leila Reddy et Rufin VanRullen --- Neurosciences, Toulouse 3, soutenue en 2022 (Président du jury),

	\item ``Extraction d'informations spatio-temporelles du mouvement avec un réseau de neurones à spike : application sur une tâche de prédiction de trajectoire de balles'' par Guillaume Debat sous la direction de Robin Baurès et de Michel Paindavoine - Toulouse 3  (Neurosciences) 	Soutenue le 17-12-2021,

	\item ``Fully event-based motion estimation of neuromorphic binocular systems'' par Laurent Dardelet sous la direction de Sio-Hoï Ieng et de Ryad Benosman - Sorbonne université (Ingénierie) Soutenue le 15-12-2021,

	% \item ``Precise timing and computationally efficient learning in neuromorphic systems'' par Omar Oubari sous la direction de Ryad Benosman et Sio-Hoï Ieng --- Ingénierie neuromorphique, Sorbonne Université, soutenue en 2020 (Président du jury).

	% \item ``Principes computationnels génériques et spécifiques à l’anticipation visuelle du mouvement'' par Selma Souihel sous la direction de Bruno Cessac - Université Côte d'Azur (Informatique) Soutenue le 18-12-2019,

	% \item ``Adaptation saccadique : un modèle d’apprentissage opérant et ses contraintes biologiques'' par Sohir Rahmouni sous la direction de Laurent Madelain - Université de Lille (2018-2021) Soutenue le 13-12-2019,

	% \item ``Apport des dispositifs de restauration de la vision et de la résolution temporelle''  par Mylène Poujade sous la direction de Ryad Benosman - Sorbonne université	(Neurosciences) 	Soutenue le 25-09-2019,

	% \item ``Implémentation multi-niveaux de modèles hiérarchiques neuronaux'' par Kévin Géhère sous la direction de Ryad Benosman - Sorbonne université 	(Sciences de l'ingénieur) Soutenue le 14-12-2018,

	% \item ``Lateral connectivity : propagation of network belief and hallucinatory-like states in the primary visual cortex'' par Benoît Le Bec sous la direction de Yves Frégnac et de Marc Pananceau - Sorbonne université (Neuroscience cognitive) Soutenue le 03-12-2018,

	% \item ``Development and validation of stimulation strategies for the optogenetics'' par Quentin Sabatier sous la direction de Ryad Benosman - Sorbonne université (Robotique) Soutenue le 07-05-2018,

	% \item ``Event-based detection and tracking'' par David Reverter Valeiras sous la direction de Ryad Benosman et de Sio-Hoï Ieng - Paris 6	(Robotique)	Soutenue le 18-09-2017,

	% \item ``Compression d'images et de vidéos inspirée du fonctionnement de la rétine''	par Effrosyni Doutsi sous la direction de Lionel Fillatre et de Marc Antonini - Université Côte d'Azur	(Automatique, traitement du signal et des images) 	Soutenue le 22-03-2017,

	% \item ``From adaptation mechanisms with visual impairment to new technologies in vision restoration'' par Sihem Kime sous la direction de Ryad Benosman et de Xavier Clady - Paris (Ingénierie) Soutenue le 20-06-2016,

	% \item ``On salience and non-accidentalness : comparing human vision to a contrario algorithms''  par Samy Blusseau sous la direction de Jean-Michel Morel - Cachan, Ecole normale supérieure (Mathématiques) Soutenue le 22-09-2015 (Examinateur).
\end{itemize} %
	

\subsection{Participation à des activités grand public} %

Je participe ou initie de nombreuses rencontres avec le grand public (cf. \url{https://laurentperrinet.github.io/project/tout-public/}) et dans la période de référence:

\begin{itemize}

    \item Sensibilisation étudiants : \og NeuroTalk sur le thème des métiers du cerveau \fg{} (\href{https://laurentperrinet.github.io/post/2025-11-24_neurotalk/}{lien vers l'événement}).

	\item Rencontre cinémas \& sciences à l'école Air Bel : Restitution du film \og \#nofakenews \fg{} (\href{https://laurentperrinet.github.io/post/2025-09-23_belair-nofakenews/}{lien vers l'article}).

    \item Conférence immersive : \og La vision, réalité ou perception ? \fg{} (\href{https://laurentperrinet.github.io/talk/2025-06-12_explore-conference-immersive}{lien vers l'événement}).
	
	\item Participation au \og Festival EXPLORE \fg{}, organisé par Aix-Marseille Université et la DR12 du CNRS (\href{https://laurentperrinet.github.io/talk/2025-06-10_explore-cine-sciences}{lien vers l'événement}).

    \item \og Le mystère de la Joconde éclairé par les neurosciences \fg{}, \textit{Cerveau et Psycho} (\href{https://laurentperrinet.github.io/post/2024-08-25-joconde/}{lien vers l'article}).

    \item \og Chats, mouches, humains : comment la vision a évolué en de multiples facettes \fg{}, article publié dans \textit{The Conversation} (\href{https://theconversation.com/le-jeu-du-cerveau-et-du-hasard-159388}{lien vers l'article}) et disponible également \href{https://laurentperrinet.github.io/publication/perrinet-24-yeux/}{ici}.

    \item Rencontre cinémas \& sciences à la prison des Baumettes : Participation à la \og Structure d’Accompagnement à la Sortie \fg{} (\href{https://laurentperrinet.github.io/post/2024-05-23-lieux-fictifs/}{lien vers l'article}).

    \item Publication d’un article pour le catalogue de l’exposition \og Vasarely, d’un art programmatique au numérique \fg{} (17 juin au 15 octobre 2023, Espace Culturel départemental Lympia de Nice). Catalogue édité par \href{https://www.decitre.fr/livres/vasarely-9788836649587.html}{Décitre} (ISBN: 978-88-366-4958-7). Plus d’informations sur l’exposition : \href{https://www.departement06.fr/culture/vasarely-d-un-art-programmatique-au-numerique-13667.html}{site officiel}.

    \item Participation au jury des \og Rencontres Internationales Sciences Et Cinémas (RISC) \fg{} (\href{https://laurentperrinet.github.io/post/2023-12-16-risc/}{lien vers l'article}).

    \item Participation à un article de dissémination pour le magazine en ligne \textit{Sciences \& Avenir}, écrit par Alice Carliez : \og Comment notre cerveau fait-il face à l’incertitude ? \fg{} (\href{https://laurentperrinet.github.io/post/2023-07-26-epsiloon/}{lien vers l'article}).


	\item Article de dissemination sur le \href{https://laurentperrinet.github.io/publication/perrinet-21-hasard/}{hasard} dans ``The conversation'' ($6200$ lectures au 16 Février 2022).

	\item Participation à une présentation ``stand up'' des NeuroStories : conférence invitée  ``\href{https://laurentperrinet.github.io/post/2019-10-07_neurostories/}{Le temps des sens}''.

	\item Article de dissemination sur le \href{https://laurentperrinet.github.io/publication/perrinet-19-temps/}{temps dans la perception} dans ``The conversation'' ($11600$ lectures au 16 Février 2022).

	\item Participation à des activités de dissémination aux des Journées de Neurologie de Langue Française (JNLF) : conférence invitée  ``\href{https://laurentperrinet.github.io/talk/2019-04-18-jnlf/}{Des illusions aux hallucinations visuelles: une porte sur la perception}".

	\item Article de dissemination sur les \href{https://theconversation.com/illusions-et-hallucinations-visuelles-une-porte-sur-la-perception-117389}{illusions visuelles} dans ``The conversation'' ($12500$ lectures au 16 Février 2022).

\end{itemize}



\subsection{Collaboration artistique} %
%\item  Projet artistique en collaboration avec Etienne Rey (artiste plasticien à la friche Belle de Mai) :

En parallèle avec les actions grand public, je développe une collaboration active avec un artiste plasticien, \href{https://laurentperrinet.github.io/authors/etienne-rey/}{Etienne Rey} (friche Belle de Mai, Marseille, voir~\url{https://laurentperrinet.github.io/project/art-science/}). Nous avons produit dans la période de référence:

\begin{itemize}

	\item Exposition Ososphère à la Laiterie (Strasbourg) de ``\href{https://laurentperrinet.github.io/post/2026-02-20_ososphere/}{Variable Density, série Delaunay}'' avec Étienne Rey.
	
	\item ``\href{https://laurentperrinet.github.io/talk/2026-01-19-art-and-science}{Lab Tour for Art - Perception Course}'' – En collaboration avec Étienne Rey, cette session fait partie du cours \href{https://www.risd.edu/academics/illustration/courses}{Perception en Provence} de la \href{https://www.risd.edu}{Rhode Island School of Design}, organisé par \href{https://www.instagram.com/cathuangart/}{Catherine Huang} avec \href{https://www.linkedin.com/in/ntolley/}{Nick Tolley}.	\years{2026},

	\item ``\href{https://laurentperrinet.github.io/post/2024-11-07_vibration-apparences/}{La vibration des apparences}'' – À l’occasion de la Biennale d’Aix-en-Provence, dans le cadre de CHRONIQUES – Biennale des Imaginaires Numériques, l’association Arts Vivants présente au musée Granet, du 8 novembre 2024 au 19 janvier 2025, une exposition consacrée à l’artiste contemporain Étienne Rey. \years{2024},
	
	\item ``\href{https://laurentperrinet.github.io/publication/perrinet-23-formes-et-perception/}{Formes et perception}'' – Publication d’un article écrit pour le catalogue de l’exposition “Vasarely, d’un art programmatique au numérique” qui a eu lieu du 17 juin au 15 octobre 2023 à l’Espace Culturel départemental Lympia de Nice. Le catalogue est édité par Décitre - (ISBN: 978-88-366-4958-7). \years{2023},
	
\end{itemize}


\subsection{Enseignement} %


\begin{itemize}

   \item 2023-04-03, 2024-05-13, 2025-05-26 : Master M4NC de l’institut NeuroMod, cours \og Prospective Innovation and Research \fg{} à \href{https://neuromod.univ-cotedazur.eu/}{\textit{Neuromod}} : \textit{Artificial neural networks and machine learning applied to the understanding of biological vision} (\href{https://laurentperrinet.github.io/talk/artificial-neural-networks-and-machine-learning-applied-to-the-understanding-of-biological-vision/}{lien vers le cours}). \years{2023-2026}

    \item 2023-03-29, 2024-04-10, 2026-03-05 : Master 1 Neurosciences et Sciences Cognitives (\href{https://laurentperrinet.github.io/talk/2023-04-05-ue-neurosciences-computationnelles/}{lien vers le cours}). \years{2023-2026}

    \item 2023-05-10, 2024-04-17, 2025-03-11 : Programme doctoral \textit{NeuroSchool PhD Program in Neuroscience} (\href{https://laurentperrinet.github.io/talk/2024-04-17-phd-program-sparse-representations/}{lien vers le cours}). \years{2023-2026}
\end{itemize}


\section{Transfert technologique, relations industrielles et valorisation} %
%Vous présenterez, pour les 5 ou 10 derniers semestres :
%vos participations à des contrats de recherche : contrats ANR, contrats européens (indiquer l?intitulé long, l?acronyme ou le sigle, la durée et les partenaires financeurs), contrats industriels, autres... (préciser votre rôle, les partenaires et les montants, le thème et le contenu des travaux, leur portée et leur impact) ;
%vos participations à des projets de créations d?entreprises ;
%vos participations à des travaux donnant lieu à des dépôts de brevets ou à des développements de logiciels (préciser votre rôle, décrire le contenu des travaux, les difficultés rencontrées et surmontées, l?impact des brevets ou logiciels) ;


\subsection{Contrats et collaborations} %
Au cours de 10 derniers semestres, j'ai eu l'occasion de collaborer sur plusieurs contrats de niveau national (ANR) et international (cf~\url{https://laurentperrinet.github.io/#grants}):


\begin{itemize}
\item soit à titre d'investigateur principal :
\begin{itemize}

	\item Porteur du projet ANR \og \textit{AgileNeuRobot} \fg{} (ANR-20-CE23-0021) : \og Robots aériens agiles biomimétiques pour le vol en conditions réelles \fg{}, dans le cadre du CES : CE23 - Intelligence Artificielle. Collaboration avec Stéphane Viollet (BioRobotique, Institut des Sciences du Mouvement) et Ryad Benosman (Institut de la Vision). Budget total : 435\,k€ (\href{https://laurentperrinet.github.io/grant/anr-anr/}{lien vers le projet}).

	\item Porteur du projet \og \textit{Polychronies} \fg{}, soutenu par une aide de 100\,k€ dans le cadre de \textit{France 2030} et de l’Initiative d’Excellence d’Aix-Marseille Université (A*MIDEX, projet n° AMX-21-RID-025) (\href{https://laurentperrinet.github.io/grant/polychronies/}{lien vers le projet}).
\item soit à titre de collaborateur :

\begin{itemize}
 \item \href{https://laurentperrinet.github.io/grant/anr-shootingstar/}{ANR ShootingStar} (2021/2024) avec Frédéric Chavane,
 \item \href{https://laurentperrinet.github.io/grant/anr-aces/}{ANR ACES} (2021/2024): ``CAssignment of credit and constraints on eye movement learning'' avec Anna Montagnini,
 \item \href{https://laurentperrinet.github.io/grant/anr-predicteye/}{ANR RubinVase} (2021/2024) avec \href{https://scholar.google.com/citations?user=22p8Wc4AAAAJ}{Dario Prandi} et \href{https://scholar.google.com/citations?user=dRVPKmkAAAAJ}{Luca Calatroni}.
\end{itemize}


\end{itemize}

\item J'ai eu le statut de collaborateur sur les contrats suivants :
\begin{itemize}
	\item  \href{https://laurentperrinet.github.io/grant/emergences/}{Emergences}: ``Near-physics emerging models for embedded AI'' (2023/2027, coordination globale du projet par le CEA).
	\item  \href{https://laurentperrinet.github.io/grant/aprovis-3-d/}{APROVIS3D}: ``Aprovis3D: Event-Based Artificial Inteligence'' (2019--2023, coordination globale du projet par Jean Martinet, université de Nice).
 	\item \href{https://laurentperrinet.github.io/grant/anr-priosens/}{ANR PRIOSENS} (2021/2025): ``Modelling behavioural and neuronal data within the active inference framework'' avec Anna Montagnini,
 	 \item \href{https://laurentperrinet.github.io/grant/anr-shootingstar/}{ANR ShootingStar} (2021/2024) avec Frédéric Chavane,
	  \item \href{https://laurentperrinet.github.io/grant/anr-aces/}{ANR ACES} (2022/2026): ``Assignment of credit and constraints on eye movement learning'' avec Anna Montagnini,
	  \item \href{https://laurentperrinet.github.io/grant/anr-rubinvase/}{ANR RubinVase} (2021/2024) avec \href{https://scholar.google.com/citations?user=22p8Wc4AAAAJ}{Dario Prandi} et \href{https://scholar.google.com/citations?user=dRVPKmkAAAAJ}{Luca Calatroni}.
	  \item \href{https://laurentperrinet.github.io/grant/anr-predicteye/}{ANR PredictEye} (2018--2022) : ``Mapping and predicting trajectories for eye movements'', avec Guillaume Masson.
	\item \href{https://laurentperrinet.github.io/grant/phd-icn/}{PhD ICN} A grant from the Ph.D. program in Integrative and Clinical Neuroscience (PhD position, 2017 / 2021).

\end{itemize}
\end{itemize}

\subsection{Développements de logiciels} %

Nous développons plusieurs lignes de recherche pour appliquer nos résultats à des problèmes concrets, sous forme de logiciels \href{https://laurentperrinet.github.io/project/open-science/}{open source}:

% \subsubsection{Mouvements des yeux et mouvement} %
\begin{itemize}

	\item \href{https://github.com/invibe/AnEMo}{AnEMo} : traitement du signal pour l'analyse des mouvements des yeux~\citep{Pasturel18anemo},
	\item \href{http://github.com/NeuralEnsemble/MotionClouds}{MotionClouds}: génération de textures pour la perception du mouvement~\citep{Sanz12,Vacher15nips,Vacher16},
	% \item \href{https://github.com/laurentperrinet/LeCheapEyeTracker}{LeCheapEyeTracker} -- \url{https://github.com/laurentperrinet/CatchTheEye} : Oculomètre minimal utilisant l'apprentissage profond;
\end{itemize}



% \subsubsection{Biologically-Inspired Computer Vision} %
% \begin{itemize}
% 	\item \href{https://github.com/laurentperrinet/openRetina}{openRetina}: caméra événementielle minimale,
% %	\item  collaboration avec Gabriel Crist\'obal (CSIC, Madrid), par exemple pour la classification d'emphysèmes~\citep{Nava13} % avec "Rodrigo Nava, J. Victor Marcos, Boris Escalante-Ram\'\irez, Gabriel Crist\'obal, Laurent U Perrinet, Ra\'ul S. J. Estépar. Advances in Texture Analysis for Emphysema Classification. 8259:214--221, 2013"
% %	\item  collaboration avec I3S (Sophia-Antipolis) pour la classification d'images industrielles en suivant la méthodolgie développée dans~\citep{PerrinetBednar15},
% 	\item \href{http://github.com/bicv/SparseHebbianLearning }{SparseHebbianLearning }: apprentissage non-supervisé d'images naturelles~\citep{Perrinet10shl,Perrinet19hulk},
% 	\item \href{http://github.com/bicv/SLIP}{Simple Library for Image Processing}: techniques de traitement de l'image, utilisé notamment dans~\citep{Perrinet15bicv,Ravello16droplets,PerrinetBednar15,Perrinet15eusipco,Perrinet16EUVIP},
% 	\item \href{http://github.com/bicv/LogGabor}{LogGabor}: représentations multi-échelles des contours~\citep{Fischer07,Fischer07cv},
% 	\item \href{http://github.com/bicv/SparseEdges}{SparseEdges}: codage épars (parcimonieux) d'images naturelles~\citep{Perrinet15bicv,PerrinetBednar15},
% 	\item \href{https://github.com/laurentperrinet/Khoei_2017_PLoSCB}{MotionParticles}: prédiction dynamique par filtrage particulaire (permet de reproduire~\citep{Perrinet12pred,Khoei13jpp,KhoeiMassonPerrinet17}).
% \end{itemize}

\subsubsection{Promotion du logiciel libre} %
Je participe à différentes initiatives afin de promouvoir les pratiques du logiciel libre
\begin{itemize}
	\item écriture régulière d'un \href{https://laurentperrinet.github.io/sciblog/}{blog scientifique},
	\item référencé sur les catalogues oifficiels de chercheurs comme \href{http://orcid.org/0000-0002-9536-010X}{orcid}
	\item participation à des réseaux sociaux à des fins de dissémination comme \href{https://neuromatch.social/@laurentperrinet}{mastodon}, \href{https://bsky.app/profile/laurentperrinet.bsky.social}{BlueSky}, \href{https://www.linkedin.com/in/laurent-perrinet-1857b9/}{LinkedIn}, \href{https://stackoverflow.com/users/234547/meduz}{stackOverflow}, \href{https://www.instagram.com/laurentperrinet/}{instagram} ou \href{https://github.com/laurentperrinet}{gitHub}.
\end{itemize}




\section{Encadrement, animation et management de la recherche}

%Vous présenterez, pour les 5 ou 10 derniers semestres :
%vos responsabilités dans l?animation de programmes ou projets français, européens ou internationaux (préciser le type de programme ou de projet, son ampleur et son impact, et décrire votre rôle) ;
%vos responsabilités et vos activités de direction d?équipe ou de laboratoire (préciser le nombre de personnes) ;
%vos autres responsabilités ou activités collectives au sein du CNRS ou plus largement (management de la recherche, fonctions ou missions d?intérêt général, participation à des conseils scientifiques et autres commissions, participation à des comités de lecture, etc.) ;
%autres.


\begin{itemize}
	
	\item Co-animation du \og \textit{Computational Neuroscience Center} (\href{https://conect-int.github.io/}{CONECT}) \fg{}, incubateur au sein de l’INT pour promouvoir la neuroscience théorique et computationnelle.
	
    \item \textbf{Responsabilités ou activités collectives au CNRS} : De Janvier 2022 à Juin 2025, j'ai été nommé à la commission interdisciplinaire \og Modélisation mathématique, informatique et physique pour les sciences du vivant \fg{} (CID51) lors du dernier mandat.
	
	\item En Novembre 2024, co-oganisation de la conférence internationale \href{https://laurentperrinet.github.io/post/2024-11-08-snufa/}{SNUFA} "Spiking Neural networks as Universal Function Approximators".
\end{itemize}

\subsection{Expertise scientifique} %


Outre ces responsabilités scientifiques, je participe à l'animation scientifique sous d'autres formes. Tout d'abord pour l'évaluation de la recherche par les chercheurs en tant que membre d'un comité éditorial ou en temps que relecteur. Je développe aussi des collaborations internationales et en même temps dans la vie sociale de l'organisme:

\begin{itemize}
	\item \href{https://orcid.org/my-orcid?orcid=0000-0002-9536-010X}{Rôle de relecteur} pour 31 articles dans des journaux internationaux à comité de relecture : Cell Reports, Nature Human Behaviour, Neuroscience and Biobehavioral Reviews, PLOS ONE, iScience, PLoS CB, PLoS Biology, Vision Research.

	\item jury pour de nombreuses bourses internationales (Belgique, Pays-Bas, Australie, Canada, États-Unis).

	% \item en 2019-2020 : missions d'expertise scientifique avec Arnaud Malvache à \href{https://unistellaroptics.com/}{Unistellar}, Marseille.

	% \item 2019-2020 : missions d'expertise scientifique en collaboration avec Sid Kouider à \href{https://www.next-mind.com/}{NextMind}, Paris.
\end{itemize}

\newpage
\section{Objectifs de recherche de Juin 2025 à Juin 2030}


Mon champ de recherche examine les caractéristiques des réseaux de neurones artificiels (RNA) appliqués à des modèles de la vision humaine ou animale. Cette recherche vise principalement à mieux comprendre le fonctionnement biologique des voies visuelles dans le système nerveux central, depuis la rétine jusqu'au contrôle du mouvement des yeux. De manière interdépendante, cette recherche permet de développer des méthodes d'apprentissage machine menant à des applications dans le domaine du traitement de l'image ou du signal. Sous cet angle, les RNA ont récemment connu un fort essor au travers de ce que l'on nomme de nos jours l'Intelligence Artificielle (IA) ou l'apprentissage profond.

Dans ce contexte, mes recherches explorent les fondements théoriques et empiriques de l'adaptation neurale, en m'intéressant à la manière dont les propriétés structurelles et fonctionnelles co-évoluent pour traiter de manière optimale les régularités statistiques des environnements naturels. Un \textbf{premier axe expérimental}, mené en collaboration avec des neurophysiologistes, étudie la réponse du système biologique à des stimulations visuelles informées par la modélisation. Un \textbf{deuxième axe applicatif} utilise ces connaissances pour construire des algorithmes efficaces adaptés à des systèmes embarqués, comme des robots aériens. Un \textbf{troisième axe théorique} examine les conséquences de ces contraintes sur les modes de représentation possibles du code neural, en se concentrant notamment sur un codage temporel précis de l'information neurale.

En pratique, mon approche en tant que chercheur du CNRS est de contribuer aux avancées scientifiques sous forme de publications dans des revues scientifiques à comité de lecture (14 dans la période de référence, dont 2 en cours d'évaluation) ainsi que de les présenter dans des conférences internationales (27~participations). Ce travail s'accompagne d'un effort de formation auprès des étudiants de Master et au niveau doctoral à travers la \og NeuroSchool \fg{} à Marseille. Je me suis aussi engagé dans les responsabilités collectives en siégeant à la commission interdisciplinaire \og Modélisation mathématique, informatique et physique pour les sciences du vivant \fg{} (CID51). Enfin, je participe activement au rayonnement de ces recherches à travers la vulgarisation (notamment dans \emph{Science \& Avenir}, \emph{Cerveau \& Psycho}, \emph{Epsiloon} ou \emph{The Conversation}) et une collaboration artistique pérenne avec Étienne Rey, artiste plasticien à la Friche Belle de Mai à Marseille.

%============================================================
\section{Rapport d'activité détaillé}
%============================================================

Mon sujet de recherche porte sur la compréhension des réseaux de neurones biologiques qui sous-tendent nos capacités visuelles. En particulier, je considère que ces capacités reposent sur des aptitudes à la prédiction, acquises par des processus d'apprentissage sur des scènes comportementalement significatives. Cette compréhension est abordée à travers des approches expérimentales, de modélisation, ainsi que de formalisation théorique.

%------------------------------------------------------------
\subsection{Axe expérimental}
%------------------------------------------------------------

Le premier axe de mes activités de recherche a pour objectif de valider expérimentalement les algorithmes de modélisation et hypothèses théoriques concernant les processus prédictifs sous-tendant nos capacités cognitives. J'ai développé une classe de stimuli appelée les \emph{Motion Clouds} qui permettent de tester les capacités de la vision à détecter des mouvements à différentes vitesses, orientations ou fréquences~\citep{SanzLeon12}.

Pour examiner les processus prédictifs dans le système visuel, la thèse de Hugo Ladret s'est concentrée sur la capacité du cortex visuel primaire (V1) à représenter des informations de différentes précisions. Grâce à une collaboration avec l'Université de Montréal, nous avons développé un protocole expérimental en neurophysiologie qui a permis de tester la réponse à différents niveaux de précision. Notre étude a révélé deux populations neuronales distinctes dans V1~: l'une codant l'orientation avec une précision fine, l'autre exhibant une invariance à la précision~--- un mécanisme inédit impliquant des boucles récurrentes. Ces résultats, publiés dans \emph{Nature Communications Biology}~\citep{Ladret23}, remettent en cause les modèles classiques de traitement visuel et ouvrent des pistes pour des algorithmes robustes au bruit sensoriel.

Cette thématique est actuellement développée par Alexandre Lainé (thèse débutée en octobre 2024, co-direction avec Sophie Denève), appliquée au Marmouset avec une technologie d'enregistrement basée sur des \emph{Neuropixels}. Une première présentation affichée a été donnée lors du congrès CNS~2025 à Florence, et un article est en cours de préparation.

%------------------------------------------------------------
\subsection{Axe applicatif}
%------------------------------------------------------------

Dans l'axe applicatif, nous avons élaboré des méthodes de traitement de l'image et des vidéos s'inspirant du fonctionnement du système visuel biologique. Durant la thèse de Victor Boutin, nous avons conçu des modèles hiérarchiques intégrant deux caractéristiques essentielles des systèmes biologiques~: une régularisation pour contrôler la parcimonie de l'activité neuronale, et des interactions descendantes permettant d'intégrer une information contextuelle. Ces travaux ont conduit à des publications sur la topographie de V1~\citep{Chavane22} et à un modèle intégré de cette aire corticale~\citep{Franciosini21}.

Grâce à une collaboration avec Stéphane Viollet (Institut des Sciences du Mouvement, Marseille) et Ryad Benosman (Institut de la Vision, Paris), et à l'obtention de l'ANR \textit{AgileNeuRobot} (ANR-20-CE23-0021) dont j'étais le porteur principal, nous avons étendu nos modèles du système visuel afin de les appliquer à des systèmes embarqués avec des contraintes énergétiques strictes. Ce projet utilise des caméras événementielles similaires dans leur nature au fonctionnement de la rétine biologique. Il a permis de recruter Antoine Grimaldi, qui a développé des algorithmes permettant la reconnaissance en temps réel~\citep{Grimaldi24}, l'apprentissage de réseaux basés sur un codage temporel précis~\citep{Grimaldi23BC}, et le calcul de cartes de temps de contact~\citep{Nunes23ICCV}.

Ce même projet ANR a permis le développement de la rétinotopie fovéée dans des algorithmes classiques d'apprentissage profond, grâce au travail doctoral de Jean-Nicolas Jérémie (soutenu en octobre 2025). Les résultats montrent une invariance à certaines transformations géométriques et une meilleure représentation de la localisation des objets~\citep{Jeremie25}.

%------------------------------------------------------------
\subsection{Axe théorique}
%------------------------------------------------------------

Mon axe théorique principal repose sur l'utilisation des théories des processus prédictifs afin de mieux comprendre le système nerveux central. Dans la période de référence, j'ai développé des hypothèses autour de la représentation de l'information prédictive, notamment la représentation des probabilités sous-jacentes à ces processus.

En contraste avec l'approche classique où la fréquence de décharge encode une valeur analogique, j'ai développé une formalisation de la représentation de l'information sous forme de motifs spatio-temporels d'impulsions neuronales, basée sur la notion de \emph{polychronie}. Grâce au projet \textit{Polychronies} (A*MIDEX AMX-21-RID-025, 2022/2025), j'ai pu développer cette théorie, donnant lieu à un article de revue~\citep{Grimaldi22polychronies} et à des algorithmes neuromorphiques~\citep{Grimaldi23bc} démontrant un gain énergétique de l'ordre de 500~fois par rapport à une méthode analogique.

%------------------------------------------------------------
\subsection{Synthèse}
%------------------------------------------------------------

Les trois axes de mes recherches sont étroitement interconnectés, ce qui nécessite une forte interdisciplinarité. Mon intégration à l'INT permet de bien couvrir la composante expérimentale, et j'ai pu développer des collaborations pour les aspects de modélisation et de formalisation théorique. Nous avons récemment créé le centre de neurosciences computationnelles \href{https://conect-int.github.io/}{CONECT} à l'INT. Une perspective est de développer cette approche grâce à des collaborations internationales, notamment à travers le centre AISSAI du CNRS.

%%%%%%%%%%%%%%%%%%%%%%%%%%%%%%%%%%%%%%%%%%%%%%%%%%%%%%%%%%%%%


%============================================================
\section{Résumé de mon activité}
%============================================================

Mon champ de recherche examine les caractéristiques des réseaux de neurones artificiels (RNA) appliqués à des modèles de la vision humaine ou animale. Cette recherche vise principalement à mieux comprendre le fonctionnement biologique des voies visuelles dans le système nerveux central, depuis la rétine jusqu'au contrôle du mouvement des yeux. De manière interdépendante, cette recherche permet de développer des méthodes d'apprentissage machine menant à des applications dans le domaine du traitement de l'image ou du signal. Sous cet angle, les RNA ont récemment connu un fort essor au travers de ce que l'on nomme de nos jours l'Intelligence Artificielle (IA) ou l'apprentissage profond.

Dans ce contexte, mes recherches explorent les fondements théoriques et empiriques de l'adaptation neurale, en m'intéressant à la manière dont les propriétés structurelles et fonctionnelles co-évoluent pour traiter de manière optimale les régularités statistiques des environnements naturels. Un \textbf{premier axe expérimental}, mené en collaboration avec des neurophysiologistes, étudie la réponse du système biologique à des stimulations visuelles informées par la modélisation. Un \textbf{deuxième axe applicatif} utilise ces connaissances pour construire des algorithmes efficaces adaptés à des systèmes embarqués, comme des robots aériens. Un \textbf{troisième axe théorique} examine les conséquences de ces contraintes sur les modes de représentation possibles du code neural, en se concentrant notamment sur un codage temporel précis de l'information neurale.

En pratique, mon approche en tant que chercheur du CNRS est de contribuer aux avancées scientifiques sous forme de publications dans des revues scientifiques à comité de lecture (14 dans la période de référence, dont 2 en cours d'évaluation) ainsi que de les présenter dans des conférences internationales (27~participations). Ce travail s'accompagne d'un effort de formation auprès des étudiants de Master et au niveau doctoral à travers la \og NeuroSchool \fg{} à Marseille. Je me suis aussi engagé dans les responsabilités collectives en siégeant à la commission interdisciplinaire \og Modélisation mathématique, informatique et physique pour les sciences du vivant \fg{} (CID51). Enfin, je participe activement au rayonnement de ces recherches à travers la vulgarisation (notamment dans \emph{Science \& Avenir}, \emph{Cerveau \& Psycho}, \emph{Epsiloon} ou \emph{The Conversation}) et une collaboration artistique pérenne avec Étienne Rey, artiste plasticien à la Friche Belle de Mai à Marseille.

%============================================================
\section{Rapport d'activité détaillé}
%============================================================

Mon sujet de recherche porte sur la compréhension des réseaux de neurones biologiques qui sous-tendent nos capacités visuelles. En particulier, je considère que ces capacités reposent sur des aptitudes à la prédiction, acquises par des processus d'apprentissage sur des scènes comportementalement significatives. Cette compréhension est abordée à travers des approches expérimentales, de modélisation, ainsi que de formalisation théorique.

%------------------------------------------------------------
\subsection{Axe expérimental}
%------------------------------------------------------------

Le premier axe de mes activités de recherche a pour objectif de valider expérimentalement les algorithmes de modélisation et hypothèses théoriques concernant les processus prédictifs sous-tendant nos capacités cognitives. J'ai développé une classe de stimuli appelée les \emph{Motion Clouds} qui permettent de tester les capacités de la vision à détecter des mouvements à différentes vitesses, orientations ou fréquences \citep*{SanzLeon12}.

Pour examiner les processus prédictifs dans le système visuel, la thèse de Hugo Ladret s'est concentrée sur la capacité du cortex visuel primaire (V1) à représenter des informations de différentes précisions. Grâce à une collaboration avec l'Université de Montréal, nous avons développé un protocole expérimental en neurophysiologie qui a permis de tester la réponse à différents niveaux de précision. Notre étude a révélé deux populations neuronales distinctes dans V1~: l'une codant l'orientation avec une précision fine, l'autre exhibant une invariance à la précision~--- un mécanisme inédit impliquant des boucles récurrentes. Ces résultats, publiés dans \emph{Nature Communications Biology} \citep*{Ladret23}, remettent en cause les modèles classiques de traitement visuel et ouvrent des pistes pour des algorithmes robustes au bruit sensoriel.

Cette thématique est actuellement développée par Alexandre Lainé (thèse débutée en octobre 2024, co-direction avec Sophie Denève), appliquée au Marmouset avec une technologie d'enregistrement basée sur des \emph{Neuropixels}. Une première présentation affichée a été donnée lors du congrès CNS~2025 à Florence, et un article est en cours de préparation.

%------------------------------------------------------------
\subsection{Axe applicatif}
%------------------------------------------------------------

Dans l'axe applicatif, nous avons élaboré des méthodes de traitement de l'image et des vidéos s'inspirant du fonctionnement du système visuel biologique. Durant la thèse de Victor Boutin, nous avons conçu des modèles hiérarchiques intégrant deux caractéristiques essentielles des systèmes biologiques~: une régularisation pour contrôler la parcimonie de l'activité neuronale, et des interactions descendantes permettant d'intégrer une information contextuelle. Ces travaux ont conduit à des publications sur la topographie de V1 \citep*{Chavane22} et à un modèle intégré de cette aire corticale \citep*{Franciosini22}.

Grâce à une collaboration avec Stéphane Viollet (Institut des Sciences du Mouvement, Marseille) et Ryad Benosman (Institut de la Vision, Paris), et à l'obtention de l'ANR \textit{AgileNeuRobot} (ANR-20-CE23-0021) dont j'étais le porteur principal, nous avons étendu nos modèles du système visuel afin de les appliquer à des systèmes embarqués avec des contraintes énergétiques strictes. Ce projet utilise des caméras événementielles similaires dans leur nature au fonctionnement de la rétine biologique. Il a permis de recruter Antoine Grimaldi, qui a développé des algorithmes permettant la reconnaissance en temps réel \citep*{Grimaldi24}, l'apprentissage de réseaux basés sur un codage temporel précis \citep*{Grimaldi23bc}, et le calcul de cartes de temps de contact \citep*{Nunes23iccv}.

Ce même projet ANR a permis le développement de la rétinotopie fovéée dans des algorithmes classiques d'apprentissage profond, grâce au travail doctoral de Jean-Nicolas Jérémie (soutenu en octobre 2025). Les résultats montrent une invariance à certaines transformations géométriques et une meilleure représentation de la localisation des objets \citep*{Jeremie25}.

%------------------------------------------------------------
\subsection{Axe théorique}
%------------------------------------------------------------

Mon axe théorique principal repose sur l'utilisation des théories des processus prédictifs afin de mieux comprendre le système nerveux central. Dans la période de référence, j'ai développé des hypothèses autour de la représentation de l'information prédictive, notamment la représentation des probabilités sous-jacentes à ces processus.

En contraste avec l'approche classique où la fréquence de décharge encode une valeur analogique, j'ai développé une formalisation de la représentation de l'information sous forme de motifs spatio-temporels d'impulsions neuronales, basée sur la notion de \emph{polychronie}. Grâce au projet \textit{Polychronies} (A*MIDEX AMX-21-RID-025, 2022/2025), j'ai pu développer cette théorie, donnant lieu à un article de revue \citep*{Grimaldi22polychronies} et à des algorithmes neuromorphiques \citep*{Grimaldi23bc} démontrant un gain énergétique de l'ordre de 500~fois par rapport à une méthode analogique.

%------------------------------------------------------------
\subsection{Synthèse}
%------------------------------------------------------------

Les trois axes de mes recherches sont étroitement interconnectés, ce qui nécessite une forte interdisciplinarité. Mon intégration à l'INT permet de bien couvrir la composante expérimentale, et j'ai pu développer des collaborations pour les aspects de modélisation et de formalisation théorique. Nous avons récemment créé le centre de neurosciences computationnelles \href{https://conect-int.github.io/}{CONECT} à l'INT. Une perspective est de développer cette approche grâce à des collaborations internationales, notamment à travers le centre AISSAI du CNRS.



\newpage
\printbibliography
\end{document} %
